\section{Conclusion}
\label{sec:conclusion}

This study establishes that the variance structure of LLM evaluation is format-conditional. Across three multiple-choice benchmarks, item selection accounts for 42--49\% of total variance, while on the free-form benchmark GSM8K, model (15\%), prompt (12\%), and item (12\%) contribute comparable shares with no single dominant source. This is a qualitative dichotomy, not a gradient. Model$\times$item interaction (33.6\%) is the hidden second-largest source: models differ in which items they find difficult, not in overall ability. Our cross-benchmark D-study framework extends the single-benchmark G-theory approach of \citet{mesting2026hidden} to format-conditional comparison and actionable design budgets.

The format-conditional structure implies distinct optimization strategies. For multiple-choice benchmarks, item pool composition is the primary reliability lever; varying prompts or seeds contributes negligible variance. For free-form benchmarks, prompt design and model selection matter as much as item selection, and evaluation protocols should systematically vary both. Across formats, multi-facet designs are 8.8$\times$ more efficient than single-condition designs (72 versus 631 items for $G \geq 0.80$). BF16 and FP16 are interchangeable under binary scoring on single-GPU inference, which eliminates precision as a required control variable. The released D-study calculator enables practitioners to determine the number of items, models, and conditions needed for a target reliability before data collection.

Three extensions are planned. An 8-model expansion will strengthen the quasi-random model facet assumption by sampling a broader population of instruction-tuned architectures. Multi-GPU tensor parallelism will test whether precision separability holds under distributed numerical non-associativity. Fine-tuning variance analysis will quantify how adaptation reshapes the item difficulty profiles that drive model$\times$item interaction.
